\documentclass[10pt,a4paper]{amsart}
\usepackage[margin=19mm]{geometry}
\usepackage{amsmath,amssymb,mathtools}
\usepackage[scr=rsfs,cal=euler]{mathalfa}
\usepackage{xfrac}
\usepackage[unicode,hidelinks,bookmarks=false]{hyperref}
\newcommand{\ie}{i.e.\ }
\newcommand{\cf}{cf.\ }
\newcommand{\resp}{respectively\ }
\newcommand{\Subsection}{\subsection}
\providecommand{\nobreakdash}{\nobreak\hbox{-}}
\setlength{\emergencystretch}{2em}
\multlinegap0pt
\usepackage{bm}
\usepackage{bbm}
\usepackage{dsfont}
\usepackage{xspace}
\usepackage{cancel}
\usepackage{mathtools}
\usepackage{amsthm}
\makeatletter
\@ifundefined{theorem}{\newtheorem{theorem}{Theorem}}{}
\@ifundefined{proposition}{\newtheorem{proposition}[theorem]{Proposition}}{}
\makeatother
\title[Garding cones and positivity of curvature operators]{Garding cones and positivity of curvature operators}
\author{Teng Huang, Jiaogen Zhang}
\date{Source: arXiv:2506.15374v1 (CC BY 4.0)}
\begin{document}
\maketitle

\noindent\emph{Source and licence.} Garding cones and positivity of curvature operators. Version arXiv:2506.15374v1 (CC BY 4.0), distributed under the Creative Commons Attribution 4.0 International licence, as recorded in the arXiv licence metadata for this exact version and shown on the abstract page https://arxiv.org/abs/2506.15374v1. Full rights notice: licenses/arxiv-2506.15374-CC-BY-4.0.txt.\par\smallskip

\noindent\textit{Editorial scope.} Counted unit: the Proposition whose source label is \texttt{cone inclusion} in the article (source0.tex lines 385--390, source label \texttt{cone inclusion}), delivered together with the article's own complete original proof of it (source0.tex lines 391--454, one \texttt{proof} environment of 64 lines). No mathematics is written, repaired or re-derived: statement and proof are the article's own text line for line. Required context retained uncounted: none. Every definition, display and label of those lines is retained unchanged. Omitted from this package and not redistributed: 1--384, 455--970. Nothing inside a retained excerpt is omitted or altered: every retained body is delivered exactly as the article's lines are written, so no omission receipt of any kind accompanies this package. Citations: the retained excerpts carry no citation command at all, so no bibliography entry of the article is transcribed and no reference is redistributed. Reference adaptation: none was needed and none was made; every reference inside the delivered unit resolves inside it, so no unresolved reference or citation is produced. Printed slips kept literal: no printed word, symbol, hypothesis or number of the counted statement or of its proof was altered. Layout and class: the article's own class is \texttt{article} with packages including amsmath, amsthm, amssymb, times, enumerate, bm, xy, mathrsfs, amscd, color; the wrapper is the lane's pinned \texttt{amsart} editorial wrapper at 10pt on A4, which restates the theorem-like environments the retained text needs and adds the article's own macro definitions that the retained text needs (none), with extra wrapper packages (bm, bbm, dsfont, xspace, cancel, mathtools). The delivered numbering is the unit's own. Assets: no retained span contains a figure environment, a graphics-inclusion command, a PDF-page inclusion, a table or a TikZ picture; no image file is copied into this package, the package declares no asset dependency and no image file is redistributed with it. Licence: version arXiv:2506.15374v1 of this article is distributed under the Creative Commons Attribution 4.0 International licence, as recorded in the arXiv licence metadata for this exact version and shown on the abstract page https://arxiv.org/abs/2506.15374v1. A full rights notice is delivered at licenses/arxiv-2506.15374-CC-BY-4.0.txt. Characters: every retained excerpt is pure ASCII, so the delivered engine, XeLaTeX with its default Latin Modern text font, reports no missing character.
\begin{proposition}\label{cone inclusion}
For every $0<\epsilon<1$, we have the inclusion relation:
\begin{equation}
\overline{\Gamma}^{+}_{2}(\alpha_{\epsilon})\subset \overline{\mathcal{P}}_{m_{\epsilon}}\,.
\end{equation}
\end{proposition}

\begin{proof}
The proof consists of several steps.\medskip

\noindent\textbf{Step 1.}  Fix an arbitrary vector $V$ with $V_{\alpha_{\epsilon}}\in \overline{\Gamma}^{+}_{2}$. Through direct computation, we establish 
\begin{equation*}
\begin{split}
2\sigma_{2}(V_{\alpha_{\epsilon}})&=\big(\sigma_{1}(V_{\alpha_{\epsilon}})\big)^2-\sum_{i=1}^{N}(\mu_{i}-\alpha_{\epsilon}\overline{\mu})^2\\
&=(1-N\alpha_{\epsilon})^{2}\overline{\mu}^{2}-\sum_{i=1}^{N}\big(\mu_{i}^{2}-2\alpha_{\epsilon}\overline{\mu}\mu_{i}+\alpha_{\epsilon}^{2}\overline{\mu}^{2} \big)\\
&=\big[(1-N\alpha_{\epsilon})^{2}+2\alpha_{\epsilon}+N\alpha_{\epsilon}^{2} \big]\overline{\mu}^{2}-\sum_{i=1}^{N}\mu_{i}^{2}\,.
\end{split}
\end{equation*}
For the first term, using the definition of $\alpha_{k}$, we derive
\begin{equation*}
\begin{split}
&(1-N\alpha_{\epsilon})^{2}+2\alpha_{\epsilon}+N\alpha_{\epsilon}^{2}\\
&=N(N-1)\alpha_{\epsilon}^{2}-2(N-1)\alpha_{\epsilon}+1\\
&=N(N-1)\big[\alpha_{\epsilon}^{2}-\frac{2}{N}\alpha_{\epsilon}+\frac{1}{N(N-1)} \big]\\
&=N(N-1)\big[(\alpha_{\epsilon}-\frac{1}{N})^{2}+\frac{1}{N^{2}(N-1)} \big]\\
&=N(N-1)\big[\frac{\epsilon^{2}}{N^{2}}+\frac{1}{N^{2}(N-1)} \big]\\
&=\frac{1+(N-1)\epsilon^2}{N}\,.
\end{split}
\end{equation*}
This establishes the fundamental inequality 
\begin{equation}\label{sigma 2}
\frac{1+(N-1)\epsilon^2}{N}\overline{\mu}^{2}\geq \sum_{i=1}^{N}\mu_{i}^{2}\,.
\end{equation}

\noindent\textbf{Step 2.} We first assume $\overline{\mu}=0$.
%For simplicity, denote $\tilde{\mu}_{i}=\mu_{i}-\alpha_{\epsilon}\overline{\mu}$ for every $i=1,\cdots,N$. We then compute from \eqref{sigma 2 definition} that
%\[
%\begin{split}
%\sum_{i=1}^{N}\tilde{\mu}_{i}^{2}=&\big(\sigma_{1}(V_{\alpha_{\epsilon}})\big)^{2}-\sigma_{2}(V_{\alpha_{\epsilon}})
%=(1-N\alpha_{\epsilon})^{2}\overline{\mu}^{2}-\sigma_{2}(V_{\alpha_{\epsilon}})
%=-\sigma_{2}(V_{\alpha_{\epsilon}}).
%\end{split}
%\]
Since $\sigma_{2}(V_{\alpha_{\epsilon}})\geq 0$,  the inequality (\ref{sigma 2}) forces $\mu_{i}=0$ for all $i$, implying $V=0$. Thus $V\in \overline{\mathcal{P}}_{m_{\epsilon}}$ trivially holds. Consequently, we may assume $\overline{\mu}>0$ in subsequent arguments. 

\noindent\textbf{Step 3.} We demonstrate that any vector $V = (\mu_{1}, \cdots, \mu_{N})$ satisfying both $\overline{\mu}>0$ and \eqref{sigma 2} belongs to $\overline{\mathcal{P}}_{m_{\epsilon}}$.
\begin{proof}[Proof of Claim]

Let $m = m_{\epsilon}$ and assume without loss of generality that $\mu_{1} \leq \cdots \leq \mu_{N}$. 
 It is sufficient to prove
\begin{equation}\label{k positive'}
c_{0}:=\sum_{i=1}^{\lfloor m \rfloor}\mu_{i}+\big(m-\lfloor m \rfloor\big)\mu_{\lfloor m \rfloor+1}\geq 0\,.
\end{equation}
Decompose the right-hand side of \eqref{sigma 2} using the index $[m]$
$$T_{1}:=\sum_{i=1}^{[m]}\mu_{i}^{2}+\big(m-[m]\big)\mu_{[m+1]}^{2} \quad \textrm{ and } \quad T_{2}:=\sum_{i=[m+2]}^{N}\mu_{i}^{2}+\big([m+1]-m\big)\mu_{[m+1]}^{2}.$$
Denote vectors in $\mathbb{R}^{N - \lfloor m \rfloor}$:
$$ \alpha=(1,\cdots,1,\sqrt{[m+1]-m})\,, \quad \beta=(\mu_{[m+2]}, \cdots, \mu_{N}, \sqrt{[m+1]-m}\mu_{[m+1]}).$$ 
These satisfy
$$|\alpha|^2=N-m\,, \qquad |\beta|^2=T_{2}\,, \qquad \langle \alpha,\beta\rangle=\overline{\mu}-c_{0}.$$
Applying the Cauchy-Schwarz inequality $|\langle \alpha,\beta\rangle|^{2}\leq |\alpha|^{2}|\beta|^{2}$ yields: 
$$T_{2}\geq \frac{(\overline{\mu}-c_{0})^{2}}{N-m}=\frac{1+(N-1)\epsilon^2}{N}(\overline{\mu}-c_{0})^{2}.$$
Substituting into \eqref{sigma 2} gives:
$$\frac{1+(N-1)\epsilon^2}{N}\overline{\mu}^{2}\geq T_{1}+\frac{1+(N-1)\epsilon^2}{N}(\overline{\mu}-c_{0})^{2}.$$
Consequently, we obtain the key inequality
\begin{equation}\label{E1}
\frac{2 + 2(N-1)\epsilon^2}{N}\overline{\mu}c_{0} \geq T_{1} + \frac{1 + (N-1)\epsilon^2}{N}c_{0}^{2} \geq 0.
\end{equation}
The non-negativity of $c_{0}$ follows from $\overline{\mu} > 0$, establishing \eqref{k positive'}. 
\end{proof}
The preceding arguments complete the proof of Proposition \ref{cone inclusion}.
\end{proof}

\end{document}
